\section{From coordinate permutations to equal path endpoints}
\label{sec:groups}

We now turn the two images of Lemma~\ref{pairs:two-scale} into the
equal-endpoint paths required by Proposition~\ref{paths:criterion}.
Recall that $\mathcal E=\R^{(\N)}$, that $K$ is the free group on
the letters $\xi_z$ for $z\in\mathcal E$, and that path weights count
the squared-distance scales of their copy factors.  The two images give
paths of very different weights.  Their endpoints need not agree.
Our task is to correct those endpoints at a cost small relative to the
larger weight.

\subsection{Corrections supported on individual coordinates}

Let $M\subset K^s$ be the monoid generated by the diagonal subgroup
$\Delta K$ and the tuples $X_b=(\xi_{b_i})_{i=1}^s$ associated to
ordered congruent copies $b$ of $A$.  Thus a chosen expression for an
element of $M$ is a path using only copy factors of weight one.
Its \emph{degree} is the number of those factors; diagonal factors have
degree zero.  Degrees always refer to chosen expressions.

Suppose two paths have endpoints $D_+,D_-\in K^s$. If we can write
\[
 D_+^{-1}D_-=uv^{-1}\qquad(u,v\in M),
\]
then $D_+u=D_-v$. Appending the chosen expressions for $u$ and $v$
therefore gives equal endpoints, with added weights $\deg u$ and
$\deg v$, respectively. We will make their sum small relative to the
larger original path weight, so that correcting the endpoints retains
the small ratio of weights.

To construct these corrections, let
\[
 N=\ker\bigl(K\longrightarrow\Z,\ \xi_z\longmapsto1\bigr).
\]
For distinct labels $i,j$, write $M_{ij}$ for the image of $M$ under
projection onto coordinates $i,j$.

\begin{lemma}\label{groups:pairs}
For every $i\ne j$, the monoid $M_{ij}$ is a group containing
$N\times N$.
\end{lemma}
\begin{proof}
Put $r=\|a_i-a_j\|>0$.  Every ordered pair $z,z'\in\mathcal E$
at distance $r$ extends to an ordered congruent copy of $A$.
Indeed, send the direction of $a_j-a_i$ to $(z'-z)/r$, extend it to
an orthonormal frame for the affine span of $A$, and translate to send
$a_i$ to $z$.  The frame can be chosen inside a finite-dimensional
coordinate subspace of $\mathcal E$.

Set $x=\xi_z$ and $y=\xi_{z'}$.  Both $(x,y)$ and $(y,x)$ therefore
belong to $M_{ij}$, and
\[
 (x,y)^{-1}
   =(x^{-1},x^{-1})(y,x)(y^{-1},y^{-1})\in M_{ij}.
\]
The diagonal generators are also invertible within $M_{ij}$, so this
monoid is a group.  It contains
\[
 (x,y)(y^{-1},y^{-1})=(xy^{-1},1),
 \qquad
 (x,y)(x^{-1},x^{-1})=(1,yx^{-1}),
\]
and their conjugates by diagonal elements.

The graph on $\mathcal E$ joining points at distance $r$ is connected.
To see this, subdivide any displacement into pieces of length at most
$2r$.  The endpoints of each piece have a common point at distance $r$:
choose an appropriate perpendicular displacement from their midpoint.
An additional coordinate supplies the perpendicular direction if needed.
Differences of letters along graph edges consequently generate, after
taking their normal closure in $K$, all differences $\xi_z\xi_{z'}^{-1}$.
That normal closure is $N$, since identifying all free generators gives
the infinite cyclic group.  The two displayed types of elements and
diagonal conjugation thus give $N\times N\subset M_{ij}$.
\end{proof}

For $g\in K$, let $g^{[i]}\in K^s$ have value $g$ in coordinate $i$
and identity in every other coordinate.  Define
\[
 L_i=\{g\in K:g^{[i]}=uv^{-1}\text{ for some }u,v\in M\}.
\]
We call $uv^{-1}$ a \emph{right fraction}.  Its cost, with chosen
expressions for $u,v$, is the sum of their degrees.  For $g\in L_i$,
let $c_i(g)$ be the minimum cost of a right fraction representing
$g^{[i]}$.  The minimum exists because its possible costs form a
nonempty set of nonnegative integers.

We will only multiply right fractions when their supports permit an
explicit construction of a new denominator.  No closure assertion for
the whole set $MM^{-1}$ is needed.

\begin{lemma}\label{groups:fractions}
Each $L_i$ is a normal subgroup of $K$, and its cost satisfies
\[
 c_i(g^{-1})=c_i(g),\qquad
 c_i(kgk^{-1})=c_i(g),\qquad
 c_i(gh)\le c_i(g)+c_i(h)
\]
for $g,h\in L_i$ and $k\in K$.
Every tuple $(g_i)_{i=1}^s$ with $g_i\in L_i$ has a represented
right fraction of cost at most $\sum_i c_i(g_i)$.
Moreover,
\begin{equation}\label{groups:derived-inclusion}
 N^{(s-2)}\subset L_i\qquad(1\le i\le s),
\end{equation}
where $N^{(0)}=N$ and $N^{(r+1)}=[N^{(r)},N^{(r)}]$.
\end{lemma}
\begin{proof}
Inversion of a right fraction interchanges its numerator and denominator.
If $d=\Delta(k)$, then
\[
 (kgk^{-1})^{[i]}=(du)(dv)^{-1}
 \quad\text{whenever}\quad g^{[i]}=uv^{-1}.
\]
These operations preserve the represented cost.  Applying each operation
in both directions proves inversion and conjugation invariance.

The following multiplication rule supplies the cost estimate.  Given
any represented right fraction $uv^{-1}$ and $g\in L_i$, put
$h=v_i^{-1}gv_i$ and choose $h^{[i]}=ab^{-1}$ of cost $c_i(h)=c_i(g)$.
Then
\begin{equation}\label{groups:single-multiplication}
 uv^{-1}g^{[i]}
   =u h^{[i]}v^{-1}
   =ua b^{-1}v^{-1}
   =(ua)(vb)^{-1}.
\end{equation}
The resulting cost is at most the previous cost plus $c_i(g)$.
This proves product closure and subadditivity for $L_i$; together with
inversion and conjugation it proves normality.  Successively applying
\eqref{groups:single-multiplication} to the single-coordinate factors of
$(g_i)_i$ proves the asserted tuple cost bound.

For the derived subgroup inclusion, we also need a multiplication rule
for fractions supported on a pair of coordinates.  Suppose that
$z=cd^{-1}$ is such a fraction, supported on a pair $J$, and let
$uv^{-1}$ be any right fraction.  Since the projection of $M$ onto $J$
is a group, choose $b\in M$ such that $t=vb$ is the identity on $J$.
The tuples $t$ and $z$ commute, because their supports are disjoint.
Consequently
\begin{equation}\label{groups:pair-multiplication}
 uv^{-1}z
   =ubt^{-1}z
   =ubzt^{-1}
   =ubc d^{-1}t^{-1}
   =(ubc)(td)^{-1}.
\end{equation}
This is a right fraction of finite cost.  We require no bound on the
cost of the auxiliary $b$.

We now prove \eqref{groups:derived-inclusion} by induction on the number
of coordinates.  The proof applies to any submonoid of $K^s$ containing
the diagonal whose pair projections are groups containing $N\times N$.
For $s=2$, the whole monoid is its pair projection and contains
$N\times N$, which proves the claim.
For $s\ge3$, fix distinct indices $i,j,k$.
The monoid obtained by projecting away $j$ satisfies the same abstract
hypotheses.  By induction, every $a\in N^{(s-3)}$ has a right fraction
in that projection supported at $i$.  Lift its numerator and denominator
to $M$.  Their quotient is a right fraction $z_a$ supported on $\{i,j\}$,
with value $a$ at $i$.  Projecting away $k$ similarly gives a right
fraction $z_b$ supported on $\{i,k\}$ with any prescribed
$b\in N^{(s-3)}$ at $i$.

The four factors of $z_a z_b z_a^{-1}z_b^{-1}$ are pair-supported
right fractions; inverses remain fractions by interchanging numerator
and denominator.  Repeated use of \eqref{groups:pair-multiplication}
makes their product a right fraction.  Its only possibly nonidentity
coordinate is $i$, where its value is $[a,b]$.
Thus $[a,b]\in L_i$.  Since $L_i$ is a subgroup, it contains the subgroup
generated by all these commutators, namely $N^{(s-2)}$.
\end{proof}

Set
\begin{equation}\label{groups:cost}
 L=N^{(s-2)},\qquad c(g)=\max_{1\le i\le s}c_i(g)\quad(g\in L).
\end{equation}
The subgroup $L$ is normal in $K$.  Lemma~\ref{groups:fractions} shows
that $c$ is finite, symmetric under inversion, subadditive, and invariant
under conjugation by every element of $K$.  We have therefore identified
a subgroup of allowable endpoint corrections and a cost for each one.
We next construct long paths for which the needed corrections have
small cost relative to their degree.

\subsection{Averaging over coordinate permutations}

Let $H$ be a finite group of coordinate permutations acting on
$\mathcal E$, and let $Z\subset\mathcal E$ be a finite $H$-invariant set.
Consider represented tuples $w=(w_x)_{x\in Z}\in K^Z$ obtained by
multiplying diagonal factors and factors
\[
 (\xi_{h x})_{x\in Z},\qquad h\in H.
\]
The degree of such an expression is the number of factors of the latter
kind.  Right multiplication by a diagonal element has degree zero.
Replacing $w_x$ by $w_{h x}$ preserves the degree: it sends a generator
indexed by $g\in H$ to the generator indexed by $gh$.

\begin{lemma}\label{groups:averaging}
Let $H$ be a finite group of coordinate permutations of $\mathcal E$,
and let $Z\subset\mathcal E$ be a finite $H$-invariant set. For every
$\delta>0$, there is a tuple $w\in K^Z$ represented by diagonal
factors and factors $(\xi_{h x})_{x\in Z}$, $h\in H$, with degree
$m>0$, for which
\begin{equation}\label{groups:orbit-bound}
 w_xw_y^{-1}\in L,\qquad c(w_xw_y^{-1})\le\delta m
 \quad\text{whenever }x,y\text{ belong to the same }H\text{-orbit}.
\end{equation}
\end{lemma}
\begin{proof}
We first arrange membership in $L$, then make the costs small.
Start with $w_x=\xi_x$, represented with degree one; its differences
$w_xw_y^{-1}$ belong to $N$.

Suppose the differences on an orbit $O$ belong to $N^{(r)}$.
Choose $x_0\in O$ and right-multiply by the diagonal element
$w_{x_0}^{-1}$.  After this operation, every value on $O$ belongs to
$N^{(r)}$, and every difference on every orbit is unchanged.
Replace the resulting tuple by
\begin{equation}\label{groups:derived-average}
 w'_x=\prod_{h\in H}w_{h x},
\end{equation}
using a fixed order on $H$.  For $x\in O$, the list $(h x)_{h\in H}$
visits every point of $O$ with the same multiplicity $|H|/|O|$.
The products in \eqref{groups:derived-average} therefore have the same
image in the abelian group $N^{(r)}/N^{(r+1)}$ for every $x\in O$.
Their differences belong to $N^{(r+1)}$.

This procedure preserves any normal-subgroup condition already obtained
on another orbit.  Indeed, if $w_xw_y^{-1}\in G$ there for a normal
subgroup $G\lhd K$, corresponding factors of the two products have the
same image in $K/G$.  Their products have the same image as well.
Process the finitely many orbits, increasing the derived depth on each
to $s-2$.  Every degree stays positive, and we obtain membership in $L$
for all within-orbit differences.  When $s=2$, the initial tuple already
has this property.

Fix one orbit $O$.  Right-normalize at $x_0\in O$ once more, so all
values on $O$ lie in $L$, and write
$B=\max_{x\in O}c(w_x)$.  Put $h=|H|$.
If $h=1$, all orbits are singletons and \eqref{groups:orbit-bound}
already holds.  Hence suppose $h\ge2$.
Repeatedly perform the operation
\begin{equation}\label{groups:contract}
 w'_x=\left(\prod_{g\in H}w_{g x}\right)w_{x_0}^{-1}.
\end{equation}
For each $x\in O$, one factor in the product is $a=w_{x_0}$.
Write the product as $P a Q$.  Then
\[
 w'_x=P(aQa^{-1}),\qquad
 c(w'_x)\le c(P)+c(Q)\le(h-1)B.
\]
All new values on $O$ still lie in $L$, and their represented degree is
$h$ times the preceding degree.  Thus, after $k$ repetitions starting
from degree $m_0$, the maximum difference cost divided by degree is at
most
\begin{equation}\label{groups:contraction-rate}
 \frac{2B}{m_0}\left(\frac{h-1}{h}\right)^k.
\end{equation}
Here the factor two uses
$c(w_xw_y^{-1})\le c(w_x)+c(w_y)$.
The bound tends to zero.

It remains to check that treating this orbit does not spoil a bound on
another orbit.  If $a_jb_j^{-1}\in L$ for $1\le j\le h$, conjugation
invariance and subadditivity give
\begin{equation}\label{groups:product-bound}
 c\bigl((a_1\cdots a_h)(b_1\cdots b_h)^{-1}\bigr)
 \le\sum_{j=1}^h c(a_jb_j^{-1}).
\end{equation}
For completeness, write $A_j=a_1\cdots a_j$ and
$B_j=b_1\cdots b_j$, with $A_0=B_0=1$.  The identity
\[
 A_jB_j^{-1}
 =\bigl(A_{j-1}(a_jb_j^{-1})A_{j-1}^{-1}\bigr)
   \bigl(A_{j-1}B_{j-1}^{-1}\bigr)
\]
proves \eqref{groups:product-bound} inductively.
Consequently the product in \eqref{groups:contract} multiplies the
maximum difference cost on any other orbit by at most $h$, while
multiplying the degree by exactly $h$.  The appended common right factor
does not change differences.  Their cost divided by degree cannot
increase.  We may therefore achieve the bound $\delta$ on one orbit
after another; there are only finitely many orbits.
\end{proof}

\subsection{Equal endpoints and completion of sufficiency}

We can now compare the two paths supplied by coordinate permutations.
The preceding lemma makes their endpoint discrepancies inexpensive,
and Lemma~\ref{groups:fractions} realizes all discrepancies by actual
paths of unit-scale copies.

\begin{proposition}\label{groups:endpoints}
Suppose $q>0$ and $f_i,g_i\in\R^\ell$ satisfy
\[
 \|f_i-f_j\|^2=\|a_i-a_j\|^2,\qquad
 \|g_i-g_j\|^2=q\|a_i-a_j\|^2
 \quad(1\le i,j\le s),
\]
and each $g_i$ is a coordinate permutation of $f_i$.
For every $\delta>0$ there exist two paths with the same endpoint and
weights $\lambda>0$ and $\mu\ge0$ such that
\[
 \frac{\mu}{\lambda}\le q+s\delta.
\]
\end{proposition}
\begin{proof}
Embed $\R^\ell$ into $\mathcal E$, let $H$ be its full coordinate
permutation group, acting trivially on the remaining coordinates, and
let $Z$ be the union of the $H$-orbits of the $f_i$ and $g_i$.
It is finite, and $f_i,g_i$ lie in the same orbit for each $i$.
Apply Lemma~\ref{groups:averaging} with the given $\delta$ and let $m>0$
be the degree of its represented tuple $w$.

Restricting the same expression to the labels $f_i$ gives the endpoint
$D_+=(w_{f_i})_{i=1}^s$ of a path of weight $m$.
Restricting it to the labels $g_i$ gives the endpoint
$D_-=(w_{g_i})_{i=1}^s$ of a path of weight $qm$.
Indeed, each counted factor acts by a coordinate permutation, so it
preserves the corresponding distances; diagonal factors stay diagonal.

For each $i$, put $r_i=w_{f_i}w_{g_i}^{-1}$.  By
\eqref{groups:orbit-bound}, $r_i\in L$ and $c(r_i)\le\delta m$.
The $i$th component of $D_+^{-1}D_-$ is
\[
 w_{f_i}^{-1}w_{g_i}=w_{f_i}^{-1}r_i^{-1}w_{f_i},
\]
so it too belongs to $L$ and has cost at most $\delta m$.
Lemma~\ref{groups:fractions} supplies represented $u,v\in M$ such that
\[
 D_+^{-1}D_-=uv^{-1},\qquad
 \deg u+\deg v\le s\delta m.
\]
Thus $D_+u=D_-v$.
Writing $a=\deg u$ and $b=\deg v$, the two resulting paths have weights
\[
 \lambda=m+a>0,\qquad \mu=qm+b\ge0,
 \qquad
 \frac{\mu}{\lambda}
    \le\frac{qm+b}{m}\le q+s\delta.
\]
\end{proof}

\begin{proof}[Completion of the sufficiency implication in
Theorem~\ref{thm:classification}]
Assume the matrix condition of the theorem. By
Proposition~\ref{alg:equivalence}, the tensor condition
\eqref{alg:tensor} holds.
For every integer $n\ge1$, Lemma~\ref{pairs:two-scale} supplies two
images as in Proposition~\ref{groups:endpoints} with
$0<q<1/(2n)$.  Apply that proposition with $\delta=1/(2sn)$.
It gives equal-endpoint paths whose weight ratio is strictly less than
$1/n$.  Proposition~\ref{paths:criterion} now implies that $A$ is Ramsey.
All choices are made after fixing $n$; no bound uniform in $n$ on the
degrees or the number of coordinates is required.
\end{proof}
